\section{Finite volume and finite components of the original adelic unitary group}
\label{sec:finite-volume-original}

This chapter closes the finite-volume input retained by SPR25.  It also
gives the finite arithmetic-component decomposition needed when an
initial theta function is studied at a fixed finite level.  The group,
its topology, and its rational subgroup are the original ones.  No
change of Haar normalization, rational form, or target space is implicit.
The one deep input is Borel's adelic finiteness theorem, applied below
after checking its hypotheses.  Its finite-volume and finite-class
conclusions are used separately.

Write \(F,E,\sigma,n,J\) for arbitrary data in the source predicate
\(\texttt{FiniteVolume}_2\): \(F,E\) are number fields,
\(\texttt{IsCMExt}(F,E,\sigma)\) holds, \(J\in M_n(E)\), and
\(\sigma(J)^t=J\), \(\det J\ne0\).  Put
\[
 B_A=A\otimes_F E,\qquad
 g^*=(1\otimes\sigma)(g)^t,\qquad
 G(A)=\{g\in\operatorname{GL}_n(B_A):g^*Jg=J\}.
 \tag{FV.1}
\]
The symbol \(J\) in \(B_A\) always means its entrywise image under
\(\texttt{inclR}\).  Thus \(G(A)\) is literally
\(\texttt{HA}(A,\sigma,J)\), rather than an isogenous group.  Write
\(\mathbb A=\mathbb A_F\), \(\mathbb A_f=\mathbb A_{F,f}\),
\(G_{\mathbb A}=G(\mathbb A)\), \(G_f=G(\mathbb A_f)\), and
\(G_\infty=G(F\otimes_{\mathbb Q}\mathbb R)\).  The rational subgroup
\(\Gamma\) is the diagonal image of \(G(F)\).

The case \(n=0\) is included throughout.  In that case all the groups
are the one-element group, its Haar measure has finite positive mass,
its unique fundamental domain up to null sets has that mass, and one
class representative suffices.  In the proofs involving matrix entries
we may therefore assume \(n>0\).

% BEGIN FV1
\paragraph{FV1: the exact unitary algebraic group.}
\label{claim:FV1}
The functor in (FV.1) is a linear algebraic group over \(F\), with the
same \(A\)-points as the source for every commutative \(F\)-algebra \(A\).
It admits a faithful closed matrix realization over \(F\), and its
inverse is
\[
                    g^{-1}=J^{-1}g^*J.
 \tag{FV.2}
\]

\emph{Proof.}
The source CM condition says that under every embedding \(E\to\mathbb C\),
\(\sigma\) becomes complex conjugation.  Injectivity of an embedding
gives \(\sigma^2=1\).  If \(\sigma=1\), the same condition makes a complex
embedding real, contradicting total complexity.  Since \([E:F]=2\),
the fixed field is consequently \(F\).  Choose \(a\) with
\(\delta=a-\sigma(a)\ne0\).  Then \(\sigma(\delta)=-\delta\),
\(d=\delta^2\in F^\times\), and \(1,\delta\) is an \(F\)-basis of \(E\).
For every \(A\), \(B_A=A\oplus A\delta\); in these coordinates
multiplication by \(a+b\delta\) has matrix
\[
                     \begin{pmatrix}a&bd\\ b&a\end{pmatrix}.
 \tag{FV.3}
\]
This realizes \(M_n(B_A)\) faithfully inside \(M_{2n}(A)\).
An \(A\)-linear endomorphism of \(B_A^n\) is \(B_A\)-linear precisely when
it commutes with multiplication by \(\delta\).  This is a system of
linear equations.  On that linear subspace the relation \(g^*Jg=J\)
is a finite system of polynomial equations with coefficients in \(F\).
It therefore cuts out a closed subgroup of \(\operatorname{GL}_{2n}\).

There is no extra invertibility condition lost in these equations.
Since \(J\) is invertible already over \(E\), the equation gives
\((J^{-1}g^*J)g=1\).  Over the commutative ring \(B_A\), taking
determinants makes \(\det g\) a unit, so \(g\) is invertible and (FV.2)
is its inverse.  Conversely every source \(\texttt{UJ}\) element
satisfies these same equations.  The realization and its inverse
coordinate readout are defined over \(F\); their equality with
\(\texttt{sigR}\), \(\texttt{JR}\), and \(\texttt{HR}\) follows directly
from (FV.1). \(\square\)
% END FV1

% BEGIN FV2
\paragraph{FV2: geometric splitting and connectedness.}
\label{claim:FV2}
The group of FV1 becomes \(\operatorname{GL}_n\) after extending scalars
to \(E\).  In particular it is smooth, geometrically connected, and
reductive.

\emph{Proof.}
Over \(E\), use
\[
 E\otimes_F E\longrightarrow E\times E,\qquad
 u\otimes v\longmapsto(uv,u\sigma(v)).
 \tag{FV.4}
\]
Its inverse exists because \(1,\delta\) is a basis and \(2\delta\ne0\).
The involution interchanges the two factors.  A matrix over this
algebra is a pair \((a,b)\), while \(J\) becomes \((J,\sigma(J))\).
The two unitary equations are
\[
 b^tJa=J,\qquad a^t\sigma(J)b=\sigma(J).
 \tag{FV.5}
\]
Since \(\sigma(J)=J^t\), either equation, with invertible \(a,b\),
implies the other; they give
\[
              b=\sigma(J)^{-1}a^{-t}\sigma(J).
 \tag{FV.6}
\]
Projection onto \(a\), with inverse (FV.6), is an algebraic group
isomorphism with \(\operatorname{GL}_n\).

For completeness, the geometric properties used here can be read
directly on this matrix group.  \(\operatorname{GL}_n\) is the nonempty
open set \(\det\ne0\) in affine \(n^2\)-space, so it is smooth and
irreducible.  Its unipotent radical is trivial by the following
matrix argument.  If a normal subgroup consisting of unipotent
matrices contained a nonidentity element, conjugate it into Jordan
form \(u=1+N\), where \(N\) has only zeros and Jordan superdiagonal
ones.  Its transpose \(u^t\) is conjugate to it (reverse the basis
of each Jordan block), so normality also puts \(u^t\) in the
subgroup.  But
\(\operatorname{tr}(uu^t)=n+\operatorname{tr}(NN^t)=n+r\),
where \(r>0\) is the number of those superdiagonal ones.  In
characteristic zero this is not the trace \(n\) of a unipotent
matrix, a contradiction.  Thus there is no nontrivial normal
unipotent subgroup.  This argument also covers \(n=1\).
Smoothness, geometric connectedness, and triviality of the unipotent
radical are preserved by the displayed splitting, proving the
assertions over \(F\). \(\square\)
% END FV2

% BEGIN FV3
\paragraph{FV3: every \(F\)-rational character is trivial.}
\label{claim:FV3}
For the connected algebraic group of FV1,
\[
                 X_F(G)=\operatorname{Hom}_F(G,\mathbb G_m)=0.
 \tag{FV.7}
\]
In particular this assertion is about algebraic characters defined
over \(F\), not about all continuous characters of an adelic group.

\emph{Proof.}
Over an algebraic closure, the regular characters of
\(\operatorname{GL}_n\) are exactly \(\det^m\), \(m\in\mathbb Z\).
Here is the elementary verification needed in this application.
Restriction to the diagonal torus is a Laurent monomial: the units
of the Laurent polynomial ring are a nonzero constant times such a
monomial, and the homomorphism law makes the constant one.
Conjugating by permutation matrices shows that all its exponents
coincide.  On every elementary one-parameter unipotent subgroup the
character is a unit in a polynomial ring in one variable, hence is
constant and equal to one.  Diagonal and elementary matrices generate
\(\operatorname{GL}_n\) over the algebraic closure, proving the claim.

In the splitting (FV.4), the nontrivial automorphism of \(E/F\)
interchanges the two matrix factors.  Formula (FV.6) therefore sends
the determinant character to its inverse; the Gram conjugation
contributes no determinant factor.  Extend this automorphism to an
algebraic closure.  If an \(F\)-defined character has geometric
exponent \(m\), Galois invariance requires \(m=-m\), hence \(m=0\).
The resulting geometric character is one, so the original character
is one as well.  FV2 gives \(G^\circ=G\), which is the identity
component appearing in Borel's criterion. \(\square\)
% END FV3

% BEGIN FV4
\paragraph{FV4: the original adelic topology and rational subgroup.}
\label{claim:FV4}
The identification in FV1 gives topological group isomorphisms
between the source groups \(G_{\mathbb A},G_f,G_\infty\) and the
respective adelic, finite adelic, and archimedean points of its
algebraic group in Borel's convention.  Under these isomorphisms
\(\texttt{HFinIn}\) is exactly the diagonal rational group.  Moreover
\[
 G_{\mathbb A}=G_\infty\times G_f
 \tag{FV.8}
\]
with commuting factors.  \(G_{\mathbb A}\) is locally compact,
Hausdorff, and second countable; \(\Gamma\) is countable, discrete,
and closed.

\emph{Proof.}
The finite free \(A\)-module \(A\otimes_FE\), with the source
\(\texttt{moduleTopology}\), is \(A^2\) through the basis \(1,\delta\).
Thus (FV.3) and its coordinate readout are continuous for
\(A=\mathbb A,\mathbb A_f,F\otimes_{\mathbb Q}\mathbb R\), and also
over each completion \(F_v\).  On the unitary equations the inverse
matrix is the continuous polynomial expression (FV.2).  Consequently
the topology induced from the source general linear group, which
records both a matrix and its inverse, agrees with the topology
induced by its entries.

Let \(r:G\hookrightarrow\operatorname{GL}_{2n}\) be (FV.3).
Borel's local integral subgroup is
\[
          K_v=G(F_v)\cap r^{-1}\operatorname{GL}_{2n}(\mathcal O_v).
 \tag{FV.9}
\]
It is compact and open.  At every finite place where \(d\) is integral,
(FV.3) shows that \(r(g)\) and \(r(g)^{-1}\) are integral if and only
if the two \(F_v\)-coordinates of every entry of \(g\) and \(g^{-1}\)
are integral.  Only finitely many places are excluded.  Hence a
family of local source matrices belongs to the restricted product
relative to the \(K_v\) exactly when its entries and inverse entries
belong to the original finite adelic tensor algebra.  The same
calculation identifies their basic open sets: outside finitely many
places both impose the integral conditions just described; at the
remaining places both use the ordinary local matrix topology.
This proves equality of adelic topologies, not merely a bijection
of abstract groups.  The infinite factors are ordinary finite
products.  Splitting \(\mathbb A\) into its infinite and finite
factors gives (FV.8), with all source projection maps unchanged.

The source \(\texttt{HFin}\) is the image of
\(\texttt{UJ}(\sigma,J)\) under entrywise \(E\to A\otimes_FE\).
This map is injective for our adelic algebras and is precisely the
diagonal map on (FV.3).  Passing to \(\texttt{HFinIn}\) changes only
the ambient subtype, not the image or its left multiplication action.

The matrix equations in FV1 form a closed subset of the finite
product of adelic additive spaces, since (FV.2) already ensures
invertibility.  This gives local compactness and the Hausdorff
property.  A restricted product over the countably many places of
a number field, with second countable local factors and compact
open integral factors, has a countable base; hence so does \(G_{\mathbb A}\).
The rational group is countable.  To check discreteness, require
each entry of \(r(\gamma)-1\) to be integral at every finite place
and sufficiently small at every infinite place.  A nonzero such
entry lies in \(\mathcal O_F\) and has a nonzero integral field norm;
its absolute norm is at least one.  The chosen infinite bounds make
that norm less than one, a contradiction.  Thus a neighborhood
of one contains no rational point except one.  A subgroup with
isolated identity in a Hausdorff group is closed: choose a
neighborhood \(V\) with \(VV^{-1}\) meeting it only at one; each
translate of \(V\) contains at most one subgroup point, which excludes
any new point in its closure. \(\square\)
% END FV4

% BEGIN FV5
\paragraph{FV5: the checked Borel finite-volume application.}
\label{claim:FV5}
The original group \(G_{\mathbb A}\) is unimodular, and
\(\Gamma\backslash G_{\mathbb A}\) has a nonzero finite
right-invariant Radon measure.

\emph{Proof and published hypothesis check.}
We use Armand Borel, \emph{Some finiteness properties of adele groups
over number fields}, Publications Math\'ematiques de l'IH\'ES
16 (1963), 5--30,
\url{https://www.numdam.org/item/PMIHES_1963__16__5_0.pdf}.
Lemma 5.5, printed pages 20--21, supplies unimodularity when the
identity component has no rational character over the ground field.
Theorem 5.6(i), printed page 21, supplies finite invariant volume on
the adelic quotient under the same character condition.

Here is the complete match.  Borel's ground field is the number field
\(F\); his algebraic matrix group is the closed \(F\)-group of FV1;
its identity component is the whole group by FV2; and its group of
\(F\)-defined rational characters is trivial by FV3.  The adele
group, local integral subgroups, and principal rational points in
his Sections 1.1--1.3 are the groups of FV4, with the source
topology and the source rational action.  No assertion of
anisotropy of the whole unitary group is needed.  The form may be
indefinite or isotropic, and no compactness of the adelic quotient
is asserted.

Borel states the quotient as \(G_{\mathbb A}/\Gamma\), with its
left action.  Inversion induces a homeomorphism
\[
   G_{\mathbb A}/\Gamma\longrightarrow
        \Gamma\backslash G_{\mathbb A},\qquad
        g\Gamma\longmapsto\Gamma g^{-1}.
 \tag{FV.10}
\]
Push the finite invariant measure through this homeomorphism.
It is right invariant, nonzero, finite, and Radon.  Call it \(\nu\).
Its scale will not be used to choose or alter a source Haar measure.
The consequence of Lemma 5.5 is unimodularity of the literal group,
so every left Haar measure on it is also right invariant. \(\square\)
% END FV5

% BEGIN FV6
\paragraph{FV6: an exact Borel fundamental domain.}
\label{claim:FV6}
For every locally compact second countable Hausdorff group \(H\)
and countable discrete closed subgroup \(\Lambda\), there is a
Borel set \(D_0\subset H\) meeting each left \(\Lambda\)-orbit
exactly once.  It is a fundamental domain for every Haar measure
on \(H\).  Its inverse quotient section is Borel.

\emph{Proof.}
Choose an identity neighborhood \(V\) with
\(VV^{-1}\cap\Lambda=\{1\}\).  A countable family of right
translates \(U_j=Vh_j\) covers \(H\).  No orbit meets one \(U_j\)
twice: if \(x,\lambda x\in Vh_j\), then
\(\lambda\in VV^{-1}\cap\Lambda\).
Set
\[
 S_j=\bigcup_{\lambda\in\Lambda}\lambda U_j,\qquad
 B_j=U_j\setminus\bigcup_{k<j}S_k,\qquad
 D_0=\bigcup_{j\ge0}B_j.
 \tag{FV.11}
\]
Each \(S_j\) is open and each \(B_j\) is Borel.  On every orbit
choose the least \(j\) for which it meets \(U_j\); its unique point
there is its unique point of \(D_0\).  Thus the sets
\(\lambda D_0\) partition \(H\) exactly.  This gives the source
\(\texttt{IsFundamentalDomain}\), whose null-measurability and
almost-everywhere partition requirements are weaker.

The quotient map \(q:H\to\Lambda\backslash H\) is open, and on
each \(U_j\) it is a homeomorphism onto the open set \(q(U_j)\):
it is a continuous open bijection there.  The Borel pieces
\(q(U_j)\setminus\bigcup_{k<j}q(U_k)\) partition the quotient.
Use the inverse of \(q|_{U_j}\) on the \(j\)-th piece.  The resulting
section \(s\) is Borel and has image \(D_0\).
The quotient is Hausdorff and second countable because \(\Lambda\)
is closed and the action has these local charts. \(\square\)
% END FV6

% BEGIN FV7
\paragraph{FV7: from finite quotient volume to every Haar fundamental domain.}
\label{claim:FV7}
Suppose \(H,\Lambda\) satisfy FV6, \(H\) is unimodular, and
\(\Lambda\backslash H\) has a finite nonzero right-invariant Radon
measure \(\nu\).  Then every left Haar measure \(\mu\) and every
source almost-everywhere fundamental domain \(D\) satisfy
\[
                         0<\mu(D)<\infty.
 \tag{FV.12}
\]
In particular no prescribed Haar normalization has to be changed.

\emph{Proof.}
Take the exact Borel section \(s\) of FV6.  For a nonnegative Borel
function \(f\) on \(H\), define its orbit sum
\[
 P f(q)=\sum_{\lambda\in\Lambda}f(\lambda s(q)),\qquad
              \int_H f\,d\mu_\nu=\int_{\Lambda\backslash H}Pf\,d\nu.
 \tag{FV.13}
\]
The sum is measurable and is independent of the chosen representative
of \(q\).  Applying this formula to indicators defines a Borel
measure: countable additivity follows from nonnegative monotone
convergence.  Right translation permutes the points of each orbit
and sends \(q\) to its right translate.  Right invariance of \(\nu\)
therefore gives \(\mu_\nu(Eh)=\mu_\nu(E)\) for every Borel \(E\).

This measure is locally finite.  If \(C\subset H\) is compact and an
orbit meets \(C\), fix one of its points \(\lambda_0s(q)\in C\).
Every other point \(\lambda s(q)\in C\) gives
\(\lambda\lambda_0^{-1}\in\Lambda\cap CC^{-1}\).  That intersection
is finite, since \(\Lambda\) is closed and discrete.  Thus
\[
 \mu_\nu(C)\le
       |\Lambda\cap CC^{-1}|\,\nu(\Lambda\backslash H)<\infty.
 \tag{FV.14}
\]
The measure is nonzero: its value on the exact section \(D_0\) is
\(\nu(\Lambda\backslash H)>0\).  Local finiteness on a locally compact
second countable Hausdorff space makes a Borel measure Radon
(apply the usual inner and outer approximation on a countable
relatively compact open cover).  Nonzero right invariance then
gives positivity on every nonempty open set: countably many right
translates of such a set cover \(H\).  Hence \(\mu_\nu\) is a right
Haar measure.

Unimodularity makes it a left Haar measure too.  Haar uniqueness
gives \(\mu_\nu=c\mu\), with \(0<c<\infty\).  Since the orbit
sum of \(1_{D_0}\) equals one everywhere, (FV.13) gives
\[
          c\mu(D_0)=\nu(\Lambda\backslash H)\in(0,\infty).
 \tag{FV.15}
\]
These are only the basic Haar uniqueness and nonnegative integration
facts; no quotient integration theorem or additional finite-volume
theorem is being imported.

Finally let \(D\) be any source fundamental domain, possibly only
null measurable.  Its translates partition almost everywhere, so
\(\sum_{\lambda}1_D(\lambda x)=1\) for \(\mu\)-almost every \(x\).
Null-measurable sets may be replaced by Borel sets modulo a null
set; the countable translates of that error remain null.
Partition \(H\) by the exact translates of \(D_0\), translate each
piece back using left Haar invariance, and use Tonelli:
\[
 \begin{aligned}
 \mu(D)
  &=\sum_{\lambda\in\Lambda}\mu(D\cap\lambda D_0)\\
  &=\sum_{\lambda\in\Lambda}
       \mu(\lambda^{-1}D\cap D_0)
    =\int_{D_0}\sum_{\lambda\in\Lambda}1_D(\lambda x)\,d\mu(x)
    =\mu(D_0).
 \end{aligned}
 \tag{FV.16}
\]
This proves (FV.12) for the original measure and the arbitrary
original fundamental domain. \(\square\)
% END FV7

% BEGIN FV8
\paragraph{FV8: the complete original finite-volume and finite-pair producers.}
\label{claim:FV8}
The source proposition \(\texttt{FiniteVolume}_2\) holds, with all
its universal quantifiers.  For every original \(D\) and every
\(V:\texttt{HSp}\,D\), there exist \(\mu,\mathcal F\) with the literal
\(\texttt{FinPair}(V,\mu,\mathcal F)\).  More strongly, one may first
fix any Haar measure on the original target and then obtain such
a fundamental domain for that measure.

\emph{Proof.}
Fix arbitrary \(F,E,\sigma,n,J\) satisfying the hypotheses of
\(\texttt{FiniteVolume}_2\), and then arbitrary source Haar
\(\mu\) and source fundamental domain \(\mathcal F\).
FV4 checks the topology and rational subgroup; FV5 supplies
unimodularity and the finite quotient measure; FV6--FV7 give
\(\mu(\mathcal F)<\infty\).  This is the full source predicate,
including its quantification over every Haar measure and every
fundamental domain.

Now specialize to \(F=D.\texttt{host}.F\),
\(E=D.\texttt{host}.E\), \(\sigma=D.\texttt{sp.cm.sig}\), and
\(J=V.\texttt{gram}\).  The stored CM datum supplies total reality,
total complexity, degree two, and conjugation under every complex
embedding.  The stored Hermitian space supplies
\(\sigma(J)^t=J\).  Its Gram determinant is nonzero: if \(Jv=0\)
with \(v\ne0\), the nonzero vector with frame coordinates \(v\)
is orthogonal to every frame vector, hence by linearity to every
vector, contradicting the stored nondegeneracy.  These checks prove
the exact source
\(\texttt{cmData\_isCMExt\_and\_gram\_isHermND}\) contract.

The topology and rational subgroup in this specialization are
\[
 G_{\mathbb A}=\texttt{GA}(D,V),\qquad
 \Gamma=\texttt{GQ}(V),
 \tag{FV.17}
\]
by definition, with no coordinate change of the actual matrices.
Choose any Haar measure, for example the source
\(\texttt{Measure.haar}\), and the exact Borel domain of FV6.
Its four \(\texttt{FinPair}\) fields are now:
\[
 \begin{array}{ll}
 \texttt{haar}:&\mu\text{ is the chosen Haar measure},\\
 \texttt{right}:&\mu\text{ is right invariant by FV5},\\
 \texttt{fd}:&\mathcal F\text{ is the domain from FV6},\\
 \texttt{fin}:&\mu(\mathcal F)\ne\infty\text{ by FV7}.
 \end{array}
 \tag{FV.18}
\]
FV7 also proves its positive volume, so the original
\(\texttt{quot\_volume\_}\allowbreak
  \texttt{ne\_zero\_ne\_top}\) conclusion is available
for every supplied \(\texttt{IsQuot}\) pair.  This closes both
\(\texttt{VolFin.finVol}\) and
\(\texttt{HolDict.}\allowbreak\texttt{exists\_finPair}\)
in their literal scopes. \(\square\)
% END FV8

% BEGIN FV9
\paragraph{FV9: every original open finite level has finitely many classes.}
\label{claim:FV9}
For every open subgroup \(K\subset G_f\), there is a finite set
\(S\subset G_f\) such that
\[
 G_f=\bigcup_{s\in S}G(F)_f\,sK.
 \tag{FV.19}
\]
Equivalently, every \(g\in G_f\) admits
\(g=\gamma_fsk\) with \(\gamma\in G(F)\), \(s\in S\), and \(k\in K\).
This is the entire original \(\texttt{exists\_finset\_classes}\)
statement; \(K\) need not be compact.

\emph{Proof.}
Borel's Theorem 5.1 in the same paper, printed page 19, states
finiteness of the double-coset space formed from the standard
integral adelic subgroup and the principal rational group.
Its only relevant hypotheses are a number field and an algebraic
matrix group over that field; FV1 and FV4 check both, including
the precise adelic convention.  Put
\(K_0=\prod_{v<\infty}K_v\), with \(K_v\) from (FV.9).
It is compact and open in \(G_f\); Borel's standard adelic subgroup
is \(G_\infty K_0\).  Invert the finite double-coset cover in his
convention.  Then take finite components, absorbing the infinite
component of each representative into \(G_\infty\).
The result is (FV.19) for \(K_0\).

For an arbitrary open \(K\), \(L=K_0\cap K\) is open in the compact
group \(K_0\), so \([K_0:L]<\infty\).  Indeed its disjoint open
cosets form an open cover of \(K_0\), which has a finite subcover.
Choose \(k_1,\ldots,k_a\) with \(K_0=\bigcup_j k_jL\).
Replace each \(s\) in the \(K_0\)-cover by the finitely many
\(sk_j\); since \(L\subset K\), this proves (FV.19).
The source group acting on the left is exactly
\(\texttt{HFinIn}(\mathbb A_f,\sigma,J)\), by FV4, so its requested
witness \(\gamma\) is obtained by diagonal base change of the same
rational matrix.  \(S\) is chosen after \(K\) and before \(g\).
\(\square\)
% END FV9

% BEGIN FV10
\paragraph{FV10: the exact real arithmetic components.}
\label{claim:FV10}
Let \(K\subset G_f\) now be compact and open, and let \(S\) be as
in FV9.  For each \(s\in S\), put
\[
 \Gamma_s=\{\gamma_\infty:\gamma\in G(F),\
                      \gamma_f\in sKs^{-1}\}\subset G_\infty.
 \tag{FV.20}
\]
This is an arithmetic, discrete subgroup of the real points of
\(\operatorname{Res}_{F/\mathbb Q}G\), in the actual real matrix
realization.  The real group is the product of the unitary groups
of the original Hermitian matrices at the real places.  It is
reductive and has compact center.  If a function \(\phi\) on
\(G_{\mathbb A}\) is left \(\Gamma\)-invariant and right
\(K\)-invariant, then
\[
       \phi_s(u)=\phi(su)\quad(u\in G_\infty)
 \tag{FV.21}
\]
is left \(\Gamma_s\)-invariant, and bounds for the finitely many
\(\phi_s\) bound \(\phi\) on the whole adelic group.

\emph{Proof.}
First make the rational real group explicit.  Regard the
\(2n\)-dimensional \(F\)-space in FV1 as a
\(2n[F:\mathbb Q]\)-dimensional \(\mathbb Q\)-space, and impose
commutation with the multiplication matrices of a
\(\mathbb Q\)-basis of \(F\), together with the unitary equations.
This constructs \(\operatorname{Res}_{F/\mathbb Q}G\) as a closed
rational matrix group.  Its \(\mathbb Q\)-points are literally
\(G(F)\).  An integral basis of \(\mathcal O_F\) identifies
\(\mathcal O_F^{2n}\) with a full \(\mathbb Z\)-lattice.  Thus
\[
 \Gamma_0=
 G(F)\cap r^{-1}\operatorname{GL}_{2n}(\mathcal O_F)
       =G(F)\cap K_0
 \tag{FV.22}
\]
is exactly the integral matrix subgroup in this rational
realization, where the last condition means all finite components.
The equality uses
\(F\cap\prod_{v<\infty}\mathcal O_v=\mathcal O_F\), applied to
the entries of both a matrix and its inverse.

Any two compact open subgroups of \(G_f\) have intersection of
finite index in each, by the compact-coset argument of FV9.
Apply this to \(sKs^{-1}\) and \(K_0\).  Taking their inverse
images in \(G(F)\) preserves the finite-index bounds.  Hence the
rational preimage in (FV.20) is commensurable with \(\Gamma_0\).
Projection to \(G_\infty\) is injective on rational matrices and
preserves these indices, so \(\Gamma_s\) is arithmetic in the
stated rational group.  Its discreteness is also direct: for any
compact \(C\subset G_\infty\), the rational matrices concerned
lie in the compact set \(C\times sKs^{-1}\) of \(G_{\mathbb A}\).
FV4 shows that only finitely many occur.  In particular the
identity is isolated in the real projection.

Since \(F\) is totally real, \(F\otimes_{\mathbb Q}\mathbb R\)
is the product over its real places, and \(E\) supplies one
complex factor above each.  The same coordinates in (FV.1) give
\[
 G_\infty=
 \prod_{v\mid\infty}
   \{a\in\operatorname{GL}_n(\mathbb C):
                  \overline a^{\,t}J_v a=J_v\}.
 \tag{FV.23}
\]
Every \(J_v\) is Hermitian and nonsingular.  Hermitian
diagonalization identifies its factor with \(U(p_v,q_v)\).
This is just a coordinate description: (FV.20)--(FV.21) retain
the original matrices.  Geometrically, restriction of scalars
becomes a product of the \(\operatorname{GL}_n\)'s in FV2, so
it is connected and reductive.  Its real center is
\(\prod_{v\mid\infty}U(1)\) when \(n>0\), hence compact.
To see the assertion about the center without any positivity
assumption on \(J_v\), the Lie algebra
\(\{X:X^*J_v+J_vX=0\}\) spans \(M_n(\mathbb C)\) over
\(\mathbb C\): decompose any matrix into its skew-adjoint
and \(i\) times skew-adjoint parts relative to \(J_v\).
A central group element commutes with this algebra and is
therefore scalar; its unitary equation is \(|z|=1\).
The algebraic center is the corresponding product of norm-one
tori.  In particular it has no real split torus of positive
dimension.  For reference, the rational-character assertion for
the restriction of scalars also follows directly: its geometric
characters are products of determinant powers in the factors,
and complex conjugation fixes every embedding of the totally
real field while negating each determinant exponent.

For an original \(D,V,\theta\), the actual CM-type indexing and
the actual maps \(\texttt{archM}\), \(\texttt{gramAt}\) identify
(FV.23) with the real product used in HCU14.  No arbitrary choice
of a compatible real representation enters.

It remains to check the function and multiplication order.
If \(\gamma_f=sks^{-1}\), then in the commuting decomposition
(FV.8)
\[
                     \gamma\,s\,u=s\,\gamma_\infty u\,k.
 \tag{FV.24}
\]
Left rational invariance and right \(K\)-invariance give
\(\phi_s(\gamma_\infty u)=\phi_s(u)\).
For arbitrary \(g\), FV9 gives \(g_f=\gamma_fsk\); hence
\[
           g=\gamma\,s\,(\gamma_\infty^{-1}g_\infty)\,k,
 \tag{FV.25}
\]
so \(\phi(g)=\phi_s(\gamma_\infty^{-1}g_\infty)\).
The maximum of finitely many component bounds is a bound
for \(\phi\).

Finally, an original open subgroup \(U\subset G_{\mathbb A}\)
with finite-level invariance always supplies such a compact
open \(K\): the pullback of \(U\cap G_f\) to \(G_f\) is open;
intersect it with \(K_0\) to obtain a compact open subgroup
(an open subgroup is also closed).  Its invariance follows
from that of \(U\cap G_f\).  Compactness is essential in the
arithmetic assertion (FV.20); FV9 alone allows arbitrary
open \(K\), for which that assertion is not claimed. \(\square\)
% END FV10

% BEGIN FV11
\paragraph{FV11: finite target pairs before every theta parameter.}
\label{claim:FV11}
Fix the original \(D\), the original \(\texttt{RealDat}\) and its
stored target spaces \(W_i=\texttt{Wc}(D,i)\), for
\(i\in\operatorname{Fin}(D.R.m)\).  There is a single family
\[
       (\mu_i,\mathcal F_i)_i,\qquad
       \forall i\ \texttt{FinPair}(W_i,\mu_i,\mathcal F_i),
 \tag{FV.26}
\]
chosen before every theta parameter \(t\), finite translation \(x\),
coefficient vector, or function to be extended.  This closes the
entire finite-pair gate of SPR25 and the finite-pair subargument
inside the source \(\texttt{exists\_extension}\).

\emph{Proof.}
FV8 is universal over all original \(D,V\), with no parameter,
model, character, coefficient, or theta-function assumption.
Apply it once to each stored \(W_i\), choosing for example the
source Haar measure for that target and a domain given by FV6.
There are finitely many indices, so these choices form (FV.26).
The same family may then be used for all \(t,x\).  The defining
Haar, right-invariance, domain, and finite-volume properties
depend only on the target and chosen measure/domain; none changes
when theta data are subsequently quantified.

In particular the Hilbert space and quotient measure used later
are the literal source objects
\[
 L^2\bigl(\texttt{nuQ}(W_i,\mu_i,\mathcal F_i)\bigr),\qquad
 \texttt{nuQ}(W_i,\mu_i,\mathcal F_i)
  =(\mu_i|_{\mathcal F_i})_*\texttt{autMk}.
 \tag{FV.27}
\]
SPR's regular-action, smooth-representative, and multiplicity-map
arguments, and ASB's estimates, may consume this same pair.
No equivalence to a differently normalized \(L^2\)-space is
needed.  The remaining initial boundedness statement is separate:
for each actual seed and its parameters it must still give a
bound before all \(g\in\texttt{GA}(D,W_i)\).  FV10 supplies the
finite arithmetic-component dictionary for that proof, but does
not itself prove rapid decay, central finiteness, or the seed's
boundedness.  Nor does this chapter construct the common
\(\texttt{RealDat}\), its readings, or its rational tests.
\(\square\)
% END FV11

\subsection*{Exact scope of the closure}

The published imports are Borel's Lemma 5.5 and Theorems 5.6(i)
and 5.1, on the algebraic group and adelic topology proved above.
The quotient-section and every-Haar/every-domain transport are
independent arguments in FV6--FV7.  The finite-volume result is
stronger than the one selected target-family gate: it proves the
full source \(\texttt{FiniteVolume}_2\).  The finite-class result
likewise has the entire source open-level quantifier, while the
real arithmetic conclusion retains the explicitly necessary
compact-open level.

No new threshold, new coefficient field, or parameter-dependent
replacement target is introduced.  Formal source correctness
is accepted throughout.  The accompanying checks concern source
identity, human proof dependencies, and document syntax; they
include no Lean build or replay.
